\documentclass[a4paper,10pt,twocolumn]{ltjarticle}
\usepackage{kaitoushu}

\pgfplotsset{
  myaxis/.style={
    axis lines=middle,
    xlabel={$x$}, ylabel={$y$},
    every axis x label/.style={at={(ticklabel* cs:1)},anchor=west},
    every axis y label/.style={at={(ticklabel* cs:1)},anchor=south},
    tick label style={font=\scriptsize},
    label style={font=\small},
    width=5.5cm, height=5cm,
    grid=none,
    samples=100,
  }
}

\begin{document}

\problemrange{3章：関数とグラフ — 練習問題1・B (p.87)}



\mondai{1}{次の2つの放物線の頂点が一致するとき，定数 $a$，$b$ の値を定めよ．
\[ y=x^{2}-4x+a,\quad y=bx^{2}-x+2 \]}
\kotae{
それぞれ平方完成して頂点を求める．

$y=x^{2}-4x+a=(x-2)^{2}+(a-4)$\quad 頂点 $(2,\;a-4)$

$y=bx^{2}-x+2=b\!\left(x-\dfrac{1}{2b}\right)^{2}+2-\dfrac{1}{4b}$\quad 頂点 $\left(\dfrac{1}{2b},\;2-\dfrac{1}{4b}\right)$

頂点が一致する条件：
\[
\dfrac{1}{2b}=2,\qquad 2-\dfrac{1}{4b}=a-4
\]

1式目より $b=\dfrac{1}{4}$．これを2式目に代入：
$2-1=a-4 \Rightarrow a=5$

$\therefore a=5,\;b=\dfrac{1}{4}$
}

\mondai{2}{関数 $y=ax^{2}-4ax+b\;(-1\leqq x\leqq 3)$ の最大値が $11$，最小値が $-7$ となるように，定数 $a$，$b$ の値を定めよ．ただし，$a>0$ とする．}
\kotae{
平方完成：$y=a(x-2)^{2}+(b-4a)$

$a>0$（下に凸）．頂点 $x=2$ は定義域 $[-1,3]$ 内．

\textbf{最小}は頂点で：$b-4a=-7\;\cdots$①

\textbf{最大}は端点のうち頂点から遠い方．
$x=-1$：$y=9a+(b-4a)=5a+b$\quad $x=3$：$y=a+(b-4a)=-3a+b$

$a>0$ より $5a+b>-3a+b$．よって $x=-1$ で最大：
$5a+b=11\;\cdots$②

②$-$①より $9a=18 \Rightarrow a=2$，$b=1$

$\therefore a=2,\;b=1$
}

\mondai{3}{2次不等式 $ax^{2}+4x+a+3>0$ がすべての実数 $x$ に対して成り立つように，定数 $a$ の値の範囲を定めよ．}
\kotae{
\textbf{(i) $a=0$ のとき}：$4x+3>0 \Leftrightarrow x>-\dfrac{3}{4}$．すべての $x$ で成立しないので不適．

\textbf{(ii) $a\neq0$ のとき}：すべての実数で正となる条件は
\[
a>0 \quad\text{かつ}\quad D<0
\]

$D=16-4a(a+3)=-4a^{2}-12a+16$

$D<0 \Leftrightarrow a^{2}+3a-4>0 \Leftrightarrow (a+4)(a-1)>0$

$\Leftrightarrow a<-4$ または $a>1$

$a>0$ と合わせて $a>1$

$\therefore a>1$
}

\mondai{4}{放物線 $y=4-x^{2}$ と $x$ 軸とで囲まれた部分に，図のように長方形 ABCD を内接させる．この長方形の周の長さが最大となるとき，AB の長さを求めよ．}
\kotae{
放物線と $x$ 軸の交点は $x=\pm2$．

$A(-x,0),\;B(x,0),\;C(x,4-x^{2}),\;D(-x,4-x^{2})$\;$(0<x<2)$ とおく．

$\mathrm{AB}=2x$，$\mathrm{BC}=4-x^{2}$．周の長さを $\ell$ とすると
\[
\ell=2(\mathrm{AB}+\mathrm{BC})=2(2x+4-x^{2})=-2(x-1)^{2}+10
\]

$0<x<2$ の範囲で $x=1$ のとき最大値 $10$．

このとき $\mathrm{AB}=2\cdot1=2$

$\therefore \mathrm{AB}=2$
}

\mondai{5}{関数 $y=x^{2}+2mx-4m$ の最小値 $s$ は $m$ のどのような関数になるか．また，$s$ は $m$ のどのような値に対して最大になるか．}
\kotae{
$x$ について平方完成：
\[
y=(x+m)^{2}-m^{2}-4m
\]

よって $x=-m$ のとき最小値
\[
s=-m^{2}-4m
\]
となる．これを $m$ の関数として平方完成：
\[
s=-(m+2)^{2}+4
\]

$\therefore s$ は $m$ の2次関数 $s=-m^{2}-4m$ であり，$m=-2$ のとき最大値 $4$ をとる．
}

\mondai{6}{正の定数 $a$ について，$y$ 軸上に点 $\mathrm{A}(0,a)$ をとり，放物線 $y=x^{2}$ の $x\geqq0$ の部分にある点 P の座標を $(\sqrt{y},\;y)$ とおく．線分 AP の長さの2乗を $L$ とするとき，次の問いに答えよ．ただし，$0\leqq y\leqq a$ とする．}
\kotae{
\textbf{(1)} $\mathrm{AP}^{2}=(\sqrt{y}-0)^{2}+(y-a)^{2}$

\[
L=y+(y-a)^{2}=y^{2}-(2a-1)y+a^{2}
\]

\textbf{(2)} $L$ を $y$ の2次関数と見て平方完成：
\[
L=\left(y-\dfrac{2a-1}{2}\right)^{2}+a^{2}-\dfrac{(2a-1)^{2}}{4}
=\left(y-\dfrac{2a-1}{2}\right)^{2}+\dfrac{4a-1}{4}
\]

頂点の $y$ 座標：$y^{*}=\dfrac{2a-1}{2}=a-\dfrac{1}{2}$

\textbf{(i)} $a\geqq\dfrac{1}{2}$ のとき：$0\leqq y^{*}\leqq a$ なので頂点で最小．
\[
L_{\min}=a-\dfrac{1}{4}\quad\left(y=a-\dfrac{1}{2}\right)
\]

\textbf{(ii)} $0<a<\dfrac{1}{2}$ のとき：$y^{*}<0$ で定義域外．$L$ は $[0,a]$ で単調増加．
\[
L_{\min}=a^{2}\quad(y=0)
\]
}

\mondai{7}{2次関数 $y=ax^{2}+bx+c$ の標準形を求め，この関数のグラフの軸と頂点の座標を求めよ．}
\kotae{
$a\neq0$．$x^{2}$ の係数でくくって平方完成：
\begin{align*}
y&=a\!\left(x^{2}+\dfrac{b}{a}x\right)+c\\
&=a\!\left\{\!\left(x+\dfrac{b}{2a}\right)^{\!2}-\dfrac{b^{2}}{4a^{2}}\right\}+c\\
&=a\!\left(x+\dfrac{b}{2a}\right)^{\!2}-\dfrac{b^{2}-4ac}{4a}
\end{align*}

\textbf{標準形}：$y=a\!\left(x+\dfrac{b}{2a}\right)^{\!2}-\dfrac{b^{2}-4ac}{4a}$

軸：$x=-\dfrac{b}{2a}$

頂点：$\left(-\dfrac{b}{2a},\;-\dfrac{b^{2}-4ac}{4a}\right)$
}

\end{document}
